\chapter{Risks Analysis } 

\section{Project Risks and Precautions}

During this project, some risks can be expected which could impact the project or the planning of the project. These risks should be addressed beforehand and prevented during the project.

\subsection{Physical Risks}
When limiting physical risks, it all comes down to paying attention and being careful about your work. Obviously, it was agreed that all should be careful and check each other's work to prevent injuries. That being said, we made some guidelines on how to properly prevent the most likely injuries.

\usubsubsection{\textbf{Cuts and Bruises Due to Tool Use}}
During this project, tools will be used and operations will be carried out that carry a certain degree of risk with them. Any type of tool has some risk of serious injury. Therefore, certain guidelines were made on the use of hand and power tools to try to minimize the risk of injuries.

\begin{enumerate}
    \item \textbf{Check.} Check the state of the tool and see if it is safe for use or if it requires maintenance and cannot be used at all.
    \item \textbf{Plan.} Plan what actions you will be performing with the tool and plan how it will work. See if you can reach all the areas you need to reach and if there are no obstructions to your work.
    \item \textbf{Execute.} Pay attention to your work while performing the tasks. Tools are nothing to joke with, and their use should always be taken seriously.
\end{enumerate}

\usubsubsection{\textbf{Burns}}
During the project, certain things have to be soldered. Since soldering irons get very hot, this poses a serious risk of burning. The rules for tool use still apply for soldering irons, but because of the added risk, there are some extra precautions to be taken. Mainly:

\begin{itemize}
    \item Clean your work environment and make sure you stay aware of where the soldering iron is, as well as making sure everyone in the work area is aware of the fact that the soldering iron is hot.
    \item After and in between usage, make sure to put the iron in the appropriate holder, as well as making sure the iron is turned off when not in active use.
\end{itemize}

\usubsubsection{\textbf{Electrocution}}
During the project, the speaker will be powered by electricity. Since the power is supplied by very strong power supplies, electrocution is a risk. Preventing this is once again a case of checking our work and staying aware of live sources. We also agreed to keep conductive materials and liquids, such as water or laptops, away from live wires.

\usubsubsection{\textbf{Injuries Due to Fights}}
Preventing fights is a prime example of keeping relations in check. There are many ways relations can suffer and therefore cause fights or other issues. Some known reasons for frustration are discussed in the TCC. Therefore, by agreeing to and following the TCC, fights are being prevented.

\usubsubsection{\textbf{Limited components}}
There is a risk involved in only a limited amount of components being available, and the group only getting a limited amount of them. This poses a restriction on the filters. For example, the filters are only allowed to contain two capacitors and two inductors. This means that these groups must be careful with their decisions, as they may waste components otherwise. This means that there is a risk involved with picking the right components.

\usubsubsection{\textbf{Communication}}
All groups work together closely, ensuring that all the components work as they should do. For example having the right phases, voltages and currents. If there is bad communication between sub groups, it could put the end result of the project at risk.l Therefore it is vital that all subgroups work together and communicate effectively. 


    
	

